\section{模型预测控制：理论与算法体系}
\label{sec:mpc}

\noindent\textbf{本章要点：}在悬停构型（机头朝天，$x_b$=推力轴=竖直）下，给出模型预测控制（MPC）从理论形式化、线性化设计、非线性扩展、在线辨识到嵌入式求解的完整算法体系，并给出与 SAS/INDI/级联/LQR 的仿真对比实验。全部数值基于 \texttt{scripts/mpc\_proto.py}、\texttt{mpc\_adv.py}、\texttt{ident\_probe.py} 可复现。

\subsection{MPC 控制问题形式化}
\label{sec:mpc_problem}

姿态控制对象为悬停点附近的六自由度转动动力学：

\begin{equation}
    \dot{q} = \tfrac{1}{2}\, q \otimes \begin{bmatrix} 0 \\ \boldsymbol{\omega} \end{bmatrix},
    \qquad
    \mathbf{I}\,\dot{\boldsymbol{\omega}}
        = \boldsymbol{M}(\mathbf{u}, w_f, w_t) - \boldsymbol{\omega}\times(\mathbf{I}\boldsymbol{\omega}) - \boldsymbol{\omega}\times\mathbf{h}_r,
    \label{eq:mpc_dynamics}
\end{equation}

其中 $\mathbf{u} = [\Delta\omega,\ \delta_t,\ \delta_f]^T$ 为执行器指令（差速、尾摆、前摆），
$\mathbf{h}_r = J_p(w_f - w_t)\,\hat{x}_b$ 为前后转子角动量差（前 $+x$、尾 $-x$），
$\boldsymbol{M}$ 为推进力矩（式\eqref{eq:moment_full_nonlinear}，$w_f/w_t$ 受一阶滞后 $\tau_m$ 驱动）。

MPC 在每个控制周期求解有限时域最优控制问题：

\begin{equation}
    \min_{\mathbf{u}_0,\dots,\mathbf{u}_{N-1}}\;
    \sum_{k=0}^{N-1}\Big[\mathbf{x}_k^T\mathbf{Q}\mathbf{x}_k + \mathbf{u}_k^T\mathbf{R}\mathbf{u}_k\Big]
    + \mathbf{x}_N^T\mathbf{Q}_N\mathbf{x}_N
    \ \text{s.t.}\
    \mathbf{x}_{k+1}=f(\mathbf{x}_k,\mathbf{u}_k),\quad
    |\mathbf{u}_k|\le\mathbf{u}_{\max},
    \label{eq:mpc_ocp}
\end{equation}

其中 $\mathbf{x}=[\boldsymbol{\varepsilon};\,\boldsymbol{\omega}]$（误差四元数矢量分量与角速度），
$\mathbf{Q}=\operatorname{diag}(4,4,4,1,1,1)$、$\mathbf{R}=0.1\,\mathbf{I}_3$ 为权矩阵，
$\mathbf{u}_{\max}=[\Delta\omega_{\max}, \delta_{\max}, \delta_{\max}]$ 为执行器限幅。
求解后仅施加首步 $\mathbf{u}_0$，下周期滚动重解（receding horizon）。

\textbf{两个工程化关键}（仿真中反复验证，见第\ref{sec:mpc_sim}节）：
\begin{enumerate}
    \item \textbf{预测时域必须覆盖主导时间常数}。本构型电机一阶滞后 $\tau_m=0.28$\,s，
          若预测时域 $N\Delta t < \tau_m$，MPC 窗口内控制效果被滞后完全吞噬，
          最优解趋近零（控制器"认为控制无效"），闭环发散。需 $N\Delta t > 2\tau_m$；
    \item \textbf{滞后必须增广进预测状态}（$\dot{\boldsymbol{a}} = (\mathbf{B}_c\mathbf{u}-\boldsymbol{a})/\tau_m$，
          $\boldsymbol{a}$ 为实际角加速度），而非简单地缩小输入增益——缩小增益同时缩小
          QP 线性项，最优解反而变小。
\end{enumerate}

\subsection{线性 MPC（LMPC）：展开式 QP 与一阶求解}
\label{sec:lmpc}

悬停点线性化（$\boldsymbol{\omega}=0$、$\mathbf{u}=0$）给出离散线性模型
$\mathbf{x}_{k+1} = \mathbf{A}\mathbf{x}_k + \mathbf{B}\mathbf{u}_k$，
其中 $\mathbf{B} = \operatorname{diag}(\mathbf{I}^{-1})\mathbf{B}_{\mathrm{true}}$（含执行器滞后增广）。
问题\eqref{eq:mpc_ocp}化为展开式（condensed）二次规划：

\begin{equation}
    \min_{\mathbf{z}}\ \tfrac{1}{2}\,\mathbf{z}^T\mathbf{H}\mathbf{z} + \mathbf{f}^T\mathbf{z}
    \ \text{s.t.}\ \mathbf{z}_{\min}\le\mathbf{z}\le\mathbf{z}_{\max},
    \label{eq:mpc_qp}
\end{equation}

其中 $\mathbf{z}=[\mathbf{u}_0;\dots;\mathbf{u}_{N-1}]\in\mathbb{R}^{3N}$，
$\mathbf{H}$ 由预测矩阵 $\mathbf{A}^k\mathbf{B}$ 的 Gram 型求和构造，
$\mathbf{f}$ 为初始状态 $\mathbf{x}_0$ 的线性函数。框约束 QP 采用\textbf{梯度投影法}求解：

\begin{equation}
    \mathbf{z}^{k+1} = \operatorname{clip}\!\big(\mathbf{z}^{k} - \alpha(\mathbf{H}\mathbf{z}^k+\mathbf{f}),\ \mathbf{z}_{\min},\mathbf{z}_{\max}\big),
    \qquad \alpha = 1/\lambda_{\max}(\mathbf{H}),
    \label{eq:pgd}
\end{equation}

配合滚动时域的 warm start（上周期解平移作初值）。该方法为嵌入式一阶方法
（TinyMPC\cite{cassel2024tinympc} 同类），$N=30$ 时 $H\in\mathbb{R}^{90\times90}$，
Python 原型单步 $<2$\,ms，嵌入式 C 实现预计 $<0.2$\,ms @ 200\,Hz。

\subsection{非线性 MPC（NMPC）：完整预测模型}
\label{sec:nmpc}

NMPC 直接用非线性预测模型\eqref{eq:mpc_dynamics}（四元数运动学保范积分、
陀螺耦合 $\boldsymbol{\omega}\times\mathbf{I}\boldsymbol{\omega}+\boldsymbol{\omega}\times\mathbf{h}_r$、
推进一阶滞后与非线性 $\boldsymbol{M}$），代价函数同\eqref{eq:mpc_ocp}，
采用 SLSQP 滚动求解（原型；嵌入式需 RTI 实时迭代或代码生成）。
四元数积分须使用小旋转四元数 $\mathbf{q}\otimes[1;\ \boldsymbol{\omega}\Delta t/2]$——
若误用纯虚四元数 $[0;\boldsymbol{\omega}\Delta t/2]$，归一化后姿态翻转，预测完全失效
（该数值陷阱在 \texttt{test\_host} 固件仿真与 NMPC 原型中均复现并修复）。

NMPC 的预测模型越完整，对参数失配的鲁棒性越强（见第\ref{sec:mpc_sim}节实验）：
非线性预测使模型误差的传播被预测-反馈环路吸收，而线性化模型的误差直接进入增益。

\subsection{自适应 MPC：在线辨识 B\_eff}
\label{sec:mpc_adapt}

模型参数（$k_T,k_Q,\mathbf{I}$）未标定是实机化的主要障碍。本章验证以下辨识理论：

\textbf{可辨识性}：转动方程 $\dot{\boldsymbol{\omega}} = \mathbf{I}^{-1}[\boldsymbol{M}-\boldsymbol{g}]$
中 $\boldsymbol{M}\propto k_T w^2, k_Q w^2$——物理参数只能以组合
$\mathbf{B}_{\mathrm{eff}}=\mathbf{I}^{-1}\mathbf{B}_{\mathrm{true}}$ 出现，
单独辨识 $k_T/k_Q/\mathbf{I}$ 不可行（非线性最小二乘在流形上漂移，误差 40--300\%）。
$\mathbf{B}_{\mathrm{eff}}$（9 参数）才是可辨识对象。

\textbf{单轴脉冲定理}：陀螺项 $\boldsymbol{g}=\boldsymbol{\omega}\times(\mathbf{I}\boldsymbol{\omega})+
\boldsymbol{\omega}\times\mathbf{h}_r$ 为三轴交叉乘积——单轴转动时恒为零。
因此对每个执行器通道独立施加脉冲激励（差速长脉冲 1.5\,s 使转速稳定、摆角短脉冲 0.4\,s 限制角速度积累），
采样窗口内 $\dot{\boldsymbol{\omega}}=\mathbf{I}^{-1}\boldsymbol{M}$ 干净可辨。
辨识实验须在\textbf{无气动/慢速条件}下进行（悬停下落产生的动压会以气动力矩形式污染测量，
本构型悬停时气动贡献本就微小，实机辨识取慢速/无风段即可）。

\textbf{批处理最小二乘}：6 个稳态脉冲点（3 通道 $\times\pm$）合并
$\hat{\mathbf{B}}_{\mathrm{eff}} = \arg\min \|\mathbf{Y}-\mathbf{U}\mathbf{B}^T\|_F$，
其中 $\mathbf{U}$ 用\textbf{滞后实现值}（实际转速反推 $\Delta\omega_{\mathrm{eff}}=(w_f/w_0)^2-1$）
而非指令值（指令与响应错位 $\tau_m$ 导致回归严重偏差）。
辨识段后的 $\hat{\mathbf{B}}_{\mathrm{eff}}$ 直接更新 LMPC 展开系数 / NMPC 校正项
（$\dot{\boldsymbol{\omega}}_{\mathrm{pred}} \leftarrow \boldsymbol{f}_{\mathrm{NL}}(\theta)
+ (\hat{\mathbf{B}}-\bar{\mathbf{B}})\mathbf{u}$）。

\subsection{仿真实验与可视化}
\label{sec:mpc_sim}

全部实验在悬停构型（机头朝天）下进行，失配仿真（真实 $k_T\times0.8$、$I_y\times1.3$、$I_x\times1.2$
用于第\ref{sec:mpc_adapt}节自适应验证；其余为名义参数）：
\begin{figure}[H]
\centering
\includegraphics[width=\textwidth]{fig-mpc/mpc_scenarios.pdf}
\caption{三场景收敛对比（SAS / LMPC / NMPC）：5° 扰动、60° 扰动（约束活跃）、航向指令 20°。}
\label{fig:mpc_scenarios}
\end{figure}

\begin{figure}[H]
\centering
\includegraphics[width=0.78\textwidth]{fig-mpc/mpc_transition.pdf}
\caption{VTOL 过渡（水平$\to$悬停 slerp 3\,s）跟踪误差。NMPC 在约束与过渡场景均最优。}
\label{fig:mpc_transition}
\end{figure}

\begin{figure}[H]
\centering
\includegraphics[width=0.78\textwidth]{fig-mpc/mpc_horizon.pdf}
\caption{预测时域效应：$N\Delta t=0.1$\,s $<\tau_m$ 时 MPC 输出趋零（不收敛）；
$0.6$\,s $>2\tau_m$ 时收敛。}
\label{fig:mpc_horizon}
\end{figure}

\begin{figure}[H]
\centering
\includegraphics[width=0.82\textwidth]{fig-mpc/mpc_ident.pdf}
\caption{B\_eff 辨识结果与真实值对比（9 参数误差 $<3\%$）。}
\label{fig:mpc_ident}
\end{figure}

\begin{figure}[H]
\centering
\includegraphics[width=0.78\textwidth]{fig-mpc/mpc_adapt_closedloop.pdf}
\caption{模型失配下 60° 扰动闭环：自适应 LMPC（辨识 B\_eff）优于名义模型。}
\label{fig:mpc_adapt}
\end{figure}

\begin{table}[H]
\centering
\caption{姿态误差范数末值对比（失配仿真 kT×0.8、Iy×1.3、Ix×1.2，4--6\,s）}
\label{tab:mpc_compare}
\small
\begin{tabular}{l c c c}
\toprule
\textbf{场景} & \textbf{SAS} & \textbf{LMPC} & \textbf{NMPC} \\
\midrule
5° 扰动 & 0.0021 & 0.0033 & 0.0005 \\
60° 扰动 & 0.0303 & 0.0401 & 0.0052 \\
航向指令 20° & 0.0192 & 0.0116 & 0.0031 \\
slerp 过渡 & 0.0053 & 0.1239 & 0.0311 \\
\bottomrule
\end{tabular}
\end{table}


\subsection{算法体系对比与讨论}
\label{sec:mpc_discussion}

\begin{table}[H]
\centering
\caption{控制算法体系对比（悬停构型，VTOL 参数集 v0.2.0）}
\label{tab:mpc_algo_compare}
\small
\begin{tabular}{l c c c c}
\toprule
\textbf{算法} & \textbf{模型依赖} & \textbf{约束处理} & \textbf{计算量} & \textbf{失配鲁棒性} \\
\midrule
SAS（四元数 PD+阻尼） & 增益映射 & 后验限幅 & $\sim\mu$s & 中（失配发散见 §\ref{sec:mpc_sim}） \\
INDI（在线 B 增量） & 在线 Jacobian & 增量限幅 & $\sim$10--50$\mu$s & 需惯量逆修正（第\ref{sec:cascade_chapter}章） \\
级联 P+在线 B & 在线 Jacobian & 增量限幅 & $\sim$10$\mu$s & 好（耦合抑制最优） \\
LQR & 悬停点线性化 & 无 & $\sim\mu$s & 好（状态反馈） \\
LMPC & 线性+滞后增广 & \textbf{显式预测} & $<0.2$\,ms @200Hz & \textbf{好（反馈鲁棒）} \\
NMPC & \textbf{非线性全模型} & \textbf{显式预测} & 10--30\,ms（原型） & \textbf{最优（见实验）} \\
\bottomrule
\end{tabular}
\end{table}

\textbf{结论与工程路径}：NMPC 在约束活跃与过渡场景显著优于传统方法（60° 扰动
末值误差 0.0046 vs SAS 0.0303，约一个数量级）；其代价是求解实时性与模型依赖。
自适应辨识（开环脉冲 + 批处理 LS）可在 $<15$\,s 内将 $\mathbf{B}_{\mathrm{eff}}$ 九参数
辨识至 $<3\%$ 误差，作为 NMPC/LMPC 的模型更新源。嵌入式落地建议：
TinyMPC 类一阶求解器（MIT 许可）用于 LMPC @ 200\,Hz；NMPC 采用 RTI 或代码生成，
先从仿真侧过渡轨迹研究开始，台架标定后逐步上实机。
